% ============================================================================
%  config.tex - Datos del informe, colores, estilos y macros
% ============================================================================

% ---------- Datos del documento (editar aquí) ----------
\newcommand{\titulo}{Robot móvil diferencial con PIC16F877A y DRV8833}
\newcommand{\subtitulo}{Taller Práctico II: Control dual de motores y robótica móvil diferencial}
\newcommand{\asignatura}{Microcontroladores y Sistemas Embebidos}
\newcommand{\programa}{Programa de Ingeniería Electrónica}
\newcommand{\facultad}{Facultad de Ingeniería}
\newcommand{\universidad}{Universidad de Ibagué}
\newcommand{\docente}{\pendiente{Nombre del docente}}
\newcommand{\fecha}{Octubre de 2026}
\newcommand{\versiondoc}{0.1 (borrador)}

% ---------- Colores ----------
\definecolor{acento}{RGB}{0,61,121}       % azul institucional aprox.
\definecolor{acentoclaro}{RGB}{220,232,245}
\definecolor{gris}{RGB}{110,110,110}
\definecolor{codigofondo}{RGB}{247,247,247}
\definecolor{codigocomentario}{RGB}{0,128,0}
\definecolor{codigoclave}{RGB}{0,0,170}
\definecolor{codigocadena}{RGB}{163,21,21}
\definecolor{pendientecolor}{RGB}{200,0,0}

% ---------- Marcas de contenido pendiente ----------
% Uso: \pendiente{qué falta}. Buscar "\pendiente" para ver lo que falta.
\newcommand{\pendiente}[1]{\textcolor{pendientecolor}{\textit{[#1]}}}

% Recuadro para fotografías o capturas que aún no se han agregado.
% Uso: \figurapendiente{descripción}{etiqueta}
\newcommand{\figurapendiente}[2]{%
  \begin{figure}[H]
    \centering
    \fbox{\parbox[c][4cm][c]{0.8\textwidth}{\centering\pendiente{Agregar: #1}}}
    \caption{#1}
    \label{#2}
  \end{figure}}

% ---------- Encabezado y pie ----------
\pagestyle{fancy}
\fancyhf{}
\fancyhead[L]{\small\color{gris}\asignatura}
\fancyhead[R]{\small\color{gris}\universidad}
\fancyfoot[L]{\small\color{gris}Versión \versiondoc}
\fancyfoot[R]{\small\thepage}
\renewcommand{\headrulewidth}{0.4pt}
\setlength{\headheight}{14pt}

% ---------- Títulos de sección ----------
\titleformat{\section}{\Large\bfseries\color{acento}}{\thesection.}{0.6em}{}
\titleformat{\subsection}{\large\bfseries\color{acento}}{\thesubsection}{0.6em}{}
\titlespacing*{\section}{0pt}{18pt}{8pt}
\titlespacing*{\subsection}{0pt}{12pt}{6pt}

% ---------- Tablas ----------
\renewcommand{\arraystretch}{1.25}
\newcolumntype{L}{>{\raggedright\arraybackslash}X}
\newcolumntype{C}{>{\centering\arraybackslash}X}

% ---------- Listados de código ----------
\lstdefinestyle{c-xc8}{
  language=C,
  basicstyle=\ttfamily\footnotesize,
  backgroundcolor=\color{codigofondo},
  commentstyle=\color{codigocomentario},
  keywordstyle=\color{codigoclave}\bfseries,
  stringstyle=\color{codigocadena},
  morekeywords={__interrupt, volatile, __delay_ms},
  numbers=left,
  numberstyle=\tiny\color{gris},
  numbersep=8pt,
  frame=single,
  rulecolor=\color{gris!40},
  breaklines=true,
  breakatwhitespace=true,
  showstringspaces=false,
  tabsize=4,
  xleftmargin=1.5em,
  framexleftmargin=1.5em,
  captionpos=t
}
\lstset{style=c-xc8}
\renewcommand{\lstlistingname}{Código}
\renewcommand{\lstlistlistingname}{Índice de códigos}

% ---------- Metadatos del PDF ----------
\hypersetup{
  pdftitle={\titulo},
  pdfsubject={\subtitulo},
  pdfkeywords={PIC16F877A, DRV8833, PWM, robótica móvil, XC8}
}
